\section{Conclusion}
\label{sec:conclusion}

The 1.0.0 development establishes certificate, indexing, execution, and
intrinsic-data seams for certified HNF and SNF in \Lean{}. Transformation
results carry chosen inverses and both inverse identities; arbitrary finite
orderings are explicit data; HNF canonicality requires one fixed ordering; and
SNF canonical data remains invariant under reindexing. Constructive Euclidean
operations support deterministic integer and rational-polynomial execution,
and pivot-measure descent removes the former factorization dependency.

\code{SmithData} and \code{SmithSignature} expose canonical diagonal
information without execution traces. Determinantal ideals recover rank and
associate factors, support the canonical mathlib PID bridge, and define
presentation-level Fitting ideals. Column-oriented finite-free presentations
derive cokernel decompositions and a \code{Module.Presentation}. A pure
certificate checker and strict \code{decide_cbv} importer close the external
trust chain. Pinned process and FFI FLINT adapters share one implementation and
remain untrusted sources of raw data. Differential tests compare canonical
factors instead of nonunique transforms, and the raw benchmark profile
separates generation, native checking, kernel checking, and end-to-end import.

The public facade, certificate schema, indexing/result structures, and
presentation direction are frozen by exact compile-time tests. Independent
axiom reports show that the executable and strict certificate roots use only
\code{propext} and \code{Quot.sound}; the reproducibility profile rebuilds the
same gates from pinned sources. Post-1.0 applications keep separate APIs,
tests, benchmarks, papers, and artifacts, so one research backend does not
delay the stable core.
